\documentclass[11pt,a4paper]{report}

% Aberstwyth MMP Project Report Template for LaTeX
%
% Authors: Neil Taylor (nst@aber.ac.uk) and Dr. Hannah Dee (hmd1@aber.ac.uk)
%
% This has been adapted from the Leeds Thesis template and the
% Group Project template for Computer Science in Aberystywth University.
%
% All comments and suggestions welcome.
%
% Template designed to be used with pdflatex: it may need alteration to
% run with a different LaTeX engine.
%
% Note - this is offered as a starting point for your work. You are not
% required to use this template and can choose to create your own document
% without it.

% This template is suitable for students with an engineering-style project,
% which will be most students in the department. If your project is a research-oriented
% project, look at the alternative template.

% To build document on the unix command line, run four commands:

% pdflatex mmp-report
% bibtex mmp-report
% pdflatex mmp-report
% pdflatex mmp-report

% you will end up with mmp-report.pdf. Before submitting, add your user ID as a prefix,
% e.g. abc01-mmp-report.pdf
\usepackage{StylesAndReferences/mmp-report}

% the following packages are used for citations - You only need to include one.
%
% Use the cite package if you are using the numeric style (e.g. IEEEannot).
% Use the natbib package if you are using the author-date style (e.g. authordate2annot).
% Only use one of these and comment out or remove the other one.
\usepackage{cite}
%\usepackage{natbib}
\usepackage{pdfpages}

%%%% Title and Section Colours %%%%
% Modify these values to change the colours used for title, sections and subsections.
% Each value is a range of 0-255 in RGB colourspace.
% Idea courtesy of discussion at
% https://www.overleaf.com/learn/latex/Using_colours_in_LaTeX
% and
% https://tex.stackexchange.com/questions/75667/change-colour-on-chapter-section-headings-lyx
%
% If you prefer to have black headers, then comment out the following lines
\definecolor{mmpTitle}{RGB}{10, 85, 145}
\definecolor{mmpSection}{RGB}{10,85,155}
\definecolor{mmpSubsection}{RGB}{79,129,189}

\chapterfont{\color{mmpTitle}}  % sets colour of chapters
\sectionfont{\color{mmpSection}}  % sets colour of sections
\subsectionfont{\color{mmpSubsection}} % sets colour subsections
\subsubsectionfont{\color{mmpSubsection}} % sets colour subsections

%%%% end of Title and Section Colours %%%%


%%%% Report Type %%%%
%% comment/uncomment depending on the type of report you want to generate
\reporttype{Engineering}
%\reporttype{Research}
%%%% end of Report Type %%%%


\begin{document}

%TC:ignore

% all of the include directives below refer to tex files
% so %TC:ignore

\title{Using LiDAR models to follow pre-recorded paths in static environments}

% Your name
\author{Matt Morgan}

% Your email
\authoremail{mam107@aber.ac.uk}

\degreeschemecode{GH76} %e.g. G400
\degreeschemetitle{Artificial Intelligence and Robotics} % e.g. Computer Science
\degreetype{BSc}

\modulecode{CS39440} % i.e. CS39440, CC39440, CS39620
\moduletitle{Major Project} % i.e. Major Project or Minor Project

\date{19th April 2019} % i.e. the date of the current version of your report

\status{Draft} % Use draft until you create the release version. Then, change this to Release.
\version{0.2}

% The title and name of your supervisor.
\supervisor{Dr Fred Labrosse}

%The email for your supervisor.
\supervisoremail{ffl@aber.ac.uk}

\maketitle

%TC:endignore
 includes cover.tex - to change the content,
% edit the tex file

\pagenumbering{roman}

% This is the front page
%TC:ignore

\title{Using LiDAR models to follow pre-recorded paths in static environments}

% Your name
\author{Matt Morgan}

% Your email
\authoremail{mam107@aber.ac.uk}

\degreeschemecode{GH76} %e.g. G400
\degreeschemetitle{Artificial Intelligence and Robotics} % e.g. Computer Science
\degreetype{BSc}

\modulecode{CS39440} % i.e. CS39440, CC39440, CS39620
\moduletitle{Major Project} % i.e. Major Project or Minor Project

\date{19th April 2019} % i.e. the date of the current version of your report

\status{Draft} % Use draft until you create the release version. Then, change this to Release.
\version{0.2}

% The title and name of your supervisor.
\supervisor{Dr Fred Labrosse}

%The email for your supervisor.
\supervisoremail{ffl@aber.ac.uk}

\maketitle

%TC:endignore


% Set up page numbering
\pagestyle{empty}

% declarations of originality
\thispagestyle{empty}

%TC:ignore

%%%
%%% You must sign the declaration of originality. 
%%%
%%% You are submitting this electronically. Therefore, to sign, you 
%%% type your name and date to replace the .... characters. 
%%%
\section*{\centering Declaration of originality}

I confirm that:

\begin{itemize}
\item{This submission is my own work, except where 
clearly indicated.}

\item{I understand that there are severe penalties for Unacceptable Academic Practice, which can lead to loss of marks or even the withholding of a degree.}
 
\item{I have read the regulations on Unacceptable Academic Practice from the University's Academic Registry (AR) and the relevant sections of the current Student Handbook of the Department of Computer Science.}
 
\item{In submitting this work I understand and agree to abide by the University's regulations governing these issues.}
\end{itemize}

\vspace{2em}
Name Matt Morgan  \\

\vspace{1em}
Date 03/05/2019 \\

%%% 
%%% We would like to make a selection of final reports available to students that take 
%%% this module in future years. To enable us to do this, we require your consent. You 
%%% are not required that you do this, but if you do give your consent, then we will have 
%%% the option to select yours as one of a number of reports as examples for other 
%%% students. If you would like to give your consent, then please include the following 
%%% text and type your name and date to replace the .... characters. 
%%% 
%%% If you do not wish to give your consent, please remove this from your report. 
%%%
\vspace{1em}
\section*{\centering Consent to share this work}

By including my name below, I hereby agree to this project's report and technical work being made available to other students and academic staff of the Aberystwyth Computer Science Department.  

\vspace{2em}
Name Matt Morgan  \\

\vspace{1em}
Date 03/05/2019 \\

%TC:endignore



\thispagestyle{empty}

%TC:ignore

\section*{\centering Acknowledgements}

I am grateful to...

I'd like to thank Fred Labrosse for giving guidance during the project and for his time spent installing additional libraries on the robot.

%TC:endignore
 % Acknowledgements

\thispagestyle{empty}

%TC:ignore

\section*{\centering Abstract}

%Include an abstract for your project. This should be approximately 300 words.

%The abstract is an overview of the work you have done. Highlight the purpose of the work and the key outcomes of the work.

A van is used to transport Aberystwyth University's large robots to and
from test sites. These robots usually have to be manually driven into
and out of the van for transportation. The Robot driving itself into a van
project is an application that can be used to automate and streamline the
process of setting up large robots for field trials by loading and unloading
robots into and out of the transport van.

This project uses a forward facing LIDAR to guide the robot into and out
of the van. This LIDAR senses a specific succession of patterns when
driving into or out of the van and these patterns are used to control the
robot to drive the correct path into and out of the van, given a roughly similar starting position.

%TC:endignore
                 % Abstract

\pagenumbering{roman}
\pagestyle{fancy}
\fancyhead{}
\fancyfoot[C]{\thepage}
\renewcommand{\headrulewidth}{0 pt}
\renewcommand{\chaptermark}[1]{\markboth{#1}{}}

\tableofcontents
\newpage
\listoffigures % comment out this line if you don't have any figures / graphics
\newpage
% \listoftables % comment out this line if you don't have any tables
% \newpage

% Set up page numbering
\pagenumbering{arabic}

\setchapterheaderfooter

%TC:endignore

% include the chapters
\chapter{Background \& Objectives}

% This section should discuss your preparation for the project, including background reading, your analysis of the problem and the process or method you have followed to help structure your work.  It is likely that you will reuse part of your outline project specification, but as you write this report at the end of the project you should have more to discuss.
%
% \textbf{Note}:
%
% \begin{itemize}
%    \item All of the sections and text in this example are for illustration purposes. The main Chapters are a good starting point, but the content and actual sections that you include are likely to be different.
%
%    \item Look at the document MMP\_SO8 Project Report and Technical Work \cite{ProjectReportTechicalWork} for additional guidance.
%
% \end {itemize}

\section{Background}
%What was your background preparation for the project? What similar systems did you assess? What was your motivation and interest in this project?
\paragraph{}
When the robot drives into the van, the data the LIDAR captures will be a specific set of patterns. These patterns should be the same each time the robot drives into the van because the van itself is a static environment where not a lot can change. By recording these specific patterns and matching what the robot currently sees to the patterns, the robot should be able to be guided into the van.
\paragraph{}
There are some interesting research papers reviewing the use of sensors and techniques to get an accurate localisation measure, in particular one titled Mobile Robot Positioning {\&} Sensors and Techniques\cite{Borenstein}. In this paper they evaluate different sensors and how they can be used effectively to track a robots position and navigate accordingly. The paper mentions model matching in relation to maps (pre-existing maps and having to build the map when running) and this points to a possible solution to the problem if the map of the area is already known.
\paragraph{}
Another relevant research paper that was looked at is titled Autonomous Driving Based on Accurate Localization Using Multilayer LiDAR and Dead Reckoning\cite{Akai} and this paper is about autonomous driving using odometry with the help of multiple LIDAR's to correct the odometry. The issue they tacked is the fact that using odometry on its own leads to unreliable localisation, but combining it with the use of multiple LIDAR's can result in an accurate localisation method. While interesting to see how the two can be combined to create accurate localisation, the vehicle they use this method on is travelling a lot faster than the vehicle that will be driving into the van and therefore this level of accuracy will not be needed.
\paragraph{}
An original idea on how to tackle the problem of which way to drive was to use computer vision to detect the edges of the van and calculate the middle point in which the robot then drives to. Accurate edge detection would result in 2 specific points being detected as the edges and the difference in distance between these two points will be able to give an indication of the rotation needed for the robot to be square with the van.
\paragraph{}
More research into model matching with the use of laserscan data came up with some lecture slides from Berkeley titled Scan Matching Overview\cite{Abbeel}. These slides were extremely useful and suggested a good way to match the incoming laser data to the model would be to use an Iterative Closest Point algorithm (ICP). This would result in a translation and rotation that could be used to guide the robot in the correct direction.
\paragraph{}
The issue with using ICP is the possibility of the algorithm taking a long time to output an answer. This is because the algorithm iterates over itself many times while closing the gap between the model scan and the source scan to get a good result. More research was therefore needed to find out if ICP can be fast enough to process in real time while the robot is moving and a suitable time limit would be 100ms due to the robot moving slowly (0.1m/s). A research paper titled Efficiet Variats of the ICP Algorithm\cite{Levoy} answered if ICP could be fast enough with its section on high speed ICP variants. It suggested that through certain ways, ICP could output a result for a 3D scan of an elephant within 30ms. Due to the data from the LIDAR being a 2D point-plane it should take even less time than a 3D point-cloud and therefore ICP is capable of returning a result within the time required.
\paragraph{}
% Subjective question here hence the use of "I"
I am interested in this project because of my interest in robotics. There were several robotics related projects which I could have chosen but using the big robots that the department has seemed like the most interesting of all of these projects. I am also interested in automation and so the motivation for working on this project was to automate a repetitive task for the department and make a system that could hopefully be used on the large robots the University has for years to come.

% \begin{itemize}
%   \item Robot likely to see a specific set of patterns when driving into/out of van.
%   \item Background prep =  year 2 robotics assignment?
%   \item Motivation/interest = Interested in robotics and locomotion with a big robot seemed cool to do. It also has a practical use for the department and will likely be used for years to come.
% \end{itemize}

\section{Analysis}
% Taking into account the problem and what you learned from the background work, what was your analysis of the problem? How did your analysis help to decompose the problem into the main tasks that you would undertake? Were there alternative approaches? Why did you choose one approach compared to the alternatives?
%
% There should be a clear statement of the objectives of the work, which you will evaluate at the end of the work.
%
% In most cases, the agreed objectives or requirements will be the result of a compromise between what would ideally have been produced and what was determined to be possible in the time available. A discussion of the process of arriving at the final list is usually appropriate.
%
% As mentioned in the lectures, think about possible security issues for the project topic. Whilst these might not be relevant for all projects, do consider if there are relevant for your project. Where there are relevant security issues, discuss how they will this affect the work that you are doing. Carry forward this discussion into relevant areas for design, implementation and testing.
\paragraph{}
Looking at what was learned from the research, the most thought out and likely best way to tackle getting the direction to drive in is using model matching with ICP. This way has been used before in robotic systems and is fast enough to work in real time on a moving vehicle. A program that takes input from the LIDAR, performs ICP model matching on the data against a model and outputs the resulting transformation is therefore needed as an objective of work.
\paragraph{}
Due to the model matching only working when there is already a model to match with, a recording program is needed as an objective of work which will record the data from the LIDAR while the robot is driven manually. This introduces a configuration phase in which the robot has to be driven into the van manually at least once before it can drive the path manually.
\paragraph{}
The final part of the program is taking the transformation from the output of the ICP model matching program and using it to guide the robot along the recorded path. This is the last objective needed for a complete system.
\paragraph{}
Because this project has to run on large robots, testing it first in a simulator before trying with the real hardware is needed. Therefore, a test environment needs to be made with all models made to scale to make the environment as close to how the real world environment.

\paragraph{}
That brings the list of objectives to: The ICP Model matching program, The recording model program, The robot control program and simulated test environment with models to scale.

%% TODO: Finish this part off.

% \begin{itemize}
%   \item Need to match a model of the pattern to what the robot currently sees.
%   \item Research into model matching suggested ICP would be a good and fast way to get the direction needed to travel
%   \item Another way to do it is to detect the edges with the lidar and work out the center of the points and drive that direction.
%   \item Objectives of work include: Making the test environment in the simulator, Making sure the robot can drive into the van manually, Making a node that records laser data at specific distances, Making a node that compares the current laser data against the model using ICP, Making a node that takes the resulting transformation and uses it to drive the robot along the same path as recorded.
%   \item Needs to drive slowly into the van so if anything goes wrong, nothing breaks.
% \end{itemize}

\section{Process}
% You need to describe briefly the life cycle model or research method that you used. You do not need to write about all of the different process models that you are aware of. Focus on the process model that you have used. It is possible that you needed to adapt an existing process model to suit your project; clearly identify what you used and how you adapted it for your needs.
\paragraph{}
The methodology that was used for this project is closely related to being agile. A sprint lasts 1 week due to the regular meetings with the project supervisor. The sprint lasts from Saturday to Friday because the meeting with the project supervisor is on Friday morning and therefore Friday afternoon is spent re-evaluating what has been achieved in the last sprint and what needs to be achieved within the next sprint.
\paragraph{}
There are several tools used to keep track of tasks and issues within sprints. A general overview was made using trello which gets updated at the start of every sprint. Any issues that occurred were put into github as tracked issues and this was used to update the tasks listed on trello. The final trello overview can be seen in figure TODO.
\paragraph{}
Sprints are extremely useful when working with robots due to the ability to keep track of the original plan of things to do but also change and update the plan based on any issues or new priorities that occur.
\paragraph{}
Chapter 2 explains the original idea for the design before any sprints occurred and chapter 3 explains the final design and what needed changing as each sprint re-evaluated and iterated upon the system's design.
% \begin{itemize}
%   \item Agile sprints every week.
%   \item At the end of each sprint have a meeting with supervisor and figure out what needs to be a priority for the next sprint.
% \end{itemize}

%\addcontentsline{toc}{chapter}{Development Process}
\chapter{Original Design}

% You should concentrate on the more important aspects of the design.
% It is essential that an overview is presented before going into detail.
% As well as describing the design adopted it must also explain what other designs were considered and why they were rejected.
% The design should describe what you expected to do, and might also explain areas that you had to revise after some investigation.

% Typically, for an object-oriented design, the discussion will focus on the choice of objects and classes and the allocation of methods to classes.
% The use made of reusable components should be described and their source referenced.
% Particularly important decisions concerning data structures usually affect the architecture of a system and so should be described here.

% How much material you include on detailed design and implementation will depend very much on the nature of the project.
% It should not be padded out. Think about the significant aspects of your system.
% For example, describe the design of the user interface if it is a critical aspect of your system, or provide detail about methods and data structures that are not trivial.
% Do not spend time on long lists of trivial items and repetitive descriptions.
% If in doubt about what is appropriate, speak to your supervisor.

% You should also identify any support tools that you used.
% You should discuss your choice of implementation tools - programming language, compilers, database management system, program development environment, etc.
% Some example sub-sections may be as follows, but the specific sections are for you to define.

\section{Overall architecture}
\paragraph{}
This project had to use ROS\cite{ROS} due to the large robots using this system exclusively and therefore has to be designed with the different parts of the system as different nodes. Details on node designs and interactions are provided below. The general design is to have three separate nodes which interact with each other. These three nodes include a recording node which records the path driven by the robot, a model matching node which matches what the robot LIDAR senses in real time with the data which the recording node gathered and a driving node that takes the resulting match and turns it into driving instructions for the robot.
\paragraph{}
Due to the fact that ICP can be implemented in many ways which affect the speed and accuracy of the match, a pre-existing ICP library is to be used to make the model matching more consistent. Tests were devised and used on a few different ICP libraries to find the best and most customisable library which is explained in the testing section.
\paragraph{}
Extra helper scripts are needed to drive the robot and manually control which model is being matched with the incoming data. Other scripts such as moving the latest recording to the model file and flipping the model so the robot can drive the same path backwards can be made for convenience.

% \begin{itemize}
%   \item The idea is to record the laserscan data while manually driving the route needed and save the data as a model for driving the route automatically.
%   \item Then using this recorded data, compare the incoming laserdata with the model data to figure out the direction to drive.
%   \item Used an ICP library because I didn't want to re-invent the wheel and there are lots out there already.
% \end{itemize}

\section{Tools}
\paragraph{}
The major tools that will be used for this project include ROS\cite{ROS} for running the software on the robot, Gazebo\cite{Gazebo} for simulating the test environment, MRPT\cite{Blanco-Claraco2019} for its ICP library\cite{MRPT-ICP}, catkin\cite{CATKIN} for compilation and TODO: think of more
% \begin{itemize}
%   \item ROS
%   \item MRPT for ICP library
%   \item Compiles using catkin_make
% \end{itemize}

\section{Languages}
\paragraph{}
When using ROS, the two main languages that can be used are Python and C++. Python is easy to use and quick to get running but may take some time with processing whereas C++ is more difficult to use but is more efficient with its processing. C++ is therefore going to be used for the nodes that needed to run at maximum efficiency such as the ICP model matching and the Automatic Driving node and Python is to be used for all the helper scripts and any early prototyping.

\section{Detailed node designs and interactions}
Here is the original detailed node designed as planned from the start. These nodes and interactions can be visualised in a diagram as shown in figure TODO.
\subsection{Recording node}
\paragraph{}
The recording node first has to open a file to write the incoming data to. This file is going to be made within a data/ folder of the root directory for the project. The recording node needs to subscribe to a few different topics one being cmd\_vel and double\_ackermann movement messages that are being sent from any other nodes to control the robot. Due to the need to capture the front LIDAR data, this node also needs to subscribe to the incoming LaserScan data so it is ready to extract the distances when a snapshot is being taken.
\paragraph{}
This node should take snapshots of data at specific times which the user must specify. It should take a snapshot when the user presses a specific button and adds the extracted LaserScan distances to a new line in the file. Any number of snapshots can be taken and should be taken at critical points along the desired path. When at the end of the path the user should take one last snapshot to mark the end of the path and then close the recording node.
\paragraph{}
After recording the full path, the robot can be driven back to the original starting position and the automatic driving part of the system can be launched which will automatically move the robot along the recorded path.

% \begin{itemize}
%   \item Recording:
%   \begin{itemize}
%     \item Picks up movement Messages
%     \item When robot has driven certain distance, take the next 3 laserscan's, average them and write to csv file.
%     \item Distance calculated by using speed given from movement * time since last command.
%   \end{itemize}

\subsection{ICP compare node}
\paragraph{}
The ICP Compare node first has to load the recorded model from the file that the recording node created. When this happens, the number of lines get counted so the number of models is known and this number is published so any nodes that need this information such as the automated driving node knowing when to stop can easily get it. This node has to subscribe to the incoming LaserScan data as this will be the source data to compare against the model which was recorded. This node also subscribes to a state topic which is used to control the model snapshot that the ICP is using as its comparison data. When the node is started, the state always starts from the first state, with that being 0, and gets changed when the subscribed state topic is changed.
\paragraph{}
When the robot publishes the LaserScan data from the front LIDAR, the distance data gets extracted and used as the source to compare against the model with ICP. The ICP algorithm is implemented using MRPT's 2D ICP library which has been tested to be fast enough to process the data in real time. More information about this testing is in the Testing chapter.
% TODO: Other libraries were looked at. NEED TO ADD EARLIER SECTION ABOUT THIS.
\paragraph{}
Each time the LaserScan data gets published, ICP is used on the extracted distances and compared with the selected model's distances. The output from the ICP library should be a 3x1 transformation with an x position translating to if the robot needs to move forwards or backwards and how far, a y position translating to if the robot needs to move sideways and how far and finally an angle in which the robot needs to align itself to. These three values will get used for an ICP direction message which will be published so the automated driving node can use the result to guide the robot along the correct path.

  % \item ICP Compare:
  % \begin{itemize}
  %   \item Load from model file.
  %   \item Starts at snapshot 0.
  %   \item Gets laserscan data from robot and compares to the current snapshot needed.
  %   \item Snapshot number is controlled by other nodes. (such as automated driving)
  %   \item Uses MRPT library, other libraries were looked at.
  %   \item Output from ICP is a 3x1 transformation with an x position, a y position and an angle for rotation.
  %   \item 3x1 transformation is sent to automated driving node.
  % \end{itemize}

\subsection{Automatic driving node}
\paragraph{}
The automated driving node will subscribe to the ICP direction topic so when the directional data gets published, it can be used to guide the robot in the correct direction. This node should also control which snapshot the ICP comparing node is looking at by publishing the value on the state topic which the ICP node is subscribed to. The state should always start from 0 which is the first model and starting position for when the model got recorded. This node is also subscribed to the number of states topic which the ICP node publishes the number of lines in the file, and therefore the number of states that exist in the model. This will get used to determine whether the robot is at the final state and can stop or not.

%   \item Automated Driving:
%   \begin{itemize}
%     \item Takes 3 transformations from ICP node and averages to get rid of anomolies.
%     \item With the average worked out, the transformation gets converted into movement as follows:
%     \begin{itemize}
%       \item Reactionary based system
%       \item x translates to speed.
%       \item y translates to direction to turn if out by a lot.
%       \item angle gets added to direction if its out by a lot.
%     \end{itemize}
%     \item When x is close enough to threashold, advance snapshot by 1.
%   \end{itemize}
% \end{itemize}

\section{Messages}
\paragraph{}
There are three main messages that the system will make use of. The first is the Laserscan message from the built in ROS sensor\_msgs. This message contains a lot of different information but the only part needed for this system are the ranges in which the LIDAR can see and therefore this information gets extracted from the message and either recorded in the recording node or used as the incoming source data for the ICP node.
\paragraph{}
The second message is a custom message that was made for this system. It will be named ICP\_direction and it will contains three floats, one for the x distance, one for the y distance and one for the angle. This message will get created every time the source laserdata gets compared to the current model and is sent to the automated driving node to be turned into driving instructions.
\paragraph{}
The last message is used for controlling the movement of the robbot. There are two different messages that can be used, the main one being the standard cmd\_vel messages which should work on all robots regardless of steering method. The other message is named double\_ackermann and is used for double ackermann steering exclusively which idris uses. Because cmd\_vel messages are more general purpose and the fact that this system is meant to work on a variety of robots, cmd\_vel messages should be the main way to control robots in this system. Double\_ackermann messages can still be used, but they are the less preferred way of using the system.
\paragraph{}
The structure of all three of these messages is shown in figure TODO.

% \begin{itemize}
%   \item Laserscan Message
%   \item ICPDirectional Message
% \end{itemize}

\section{User configuration}
\paragraph{}
This system requires some set-up before the robot can drive the path automatically. First the robot needs to have a clear starting and ending position. The robot also needs to be placed at the starting position in the orientation best suited to get to the end position as the steering is for correction rather than turning.

\paragraph{}
Once the robot is in position, the user has to start the recording node and manually drive the robot along the desired path while choosing good places to take snapshots. When the robot is at the end of its path, the recording node needs to be stopped and the robot needs to be driven back to its starting location. The newly recorded path file needs to be put in the model file that the ICP node will load from and there will be a python script that gets the latest recorded file and uses it as the model.

\paragraph{}
When the robot is roughly at the same starting position, the automated node can be started and the robot will drive along the last recorded path until its completion where it will stop and display a message stating it is at the final state and close enough to stop.
% \begin{itemize}
%   \item Tried to make it easy
%   \item 1. Start recording program and move robot along desired route. Stop program when finished
%   \item 2. Place robot roughly in the same starting position so it can see roughly the same starting things
%   \item 3. Start relivant automated launch file (launch files are named after robot steering systems)
%   \item 4. Watch as the robot drives the same path which was recorded.
% \end{itemize}

\section{Additional scripts}
\paragraph{}
A few additional scripts that help debug the system or help the user automate the process will be used. These scripts include getting the most recent recording and putting the data into the model file ready for launching the automated driving system, flipping the data in the model file so the robot can use the same data to drive the route in reverse, manual control of the robot so the recording part of the system works and manual control of the model states for debugging any issues that occur if a state has been missed.
% \begin{itemize}
%   \item Get latest recording and put it in the model script
%   \item Get the model and flip the values so it can drive the same course in reverse.
%   \item Manual control of robot states
%   \item Manual control of robot
% \end{itemize}

\section{Polish}
\paragraph{}
Because this system is designed to work on all robots, the steering needs to be changed depending on the steering method used by each robot. Therefore, as an extra addition, there needs to be different launch files which contain a parameter for each of the different steering methods. This parameter changes how the program steers when reversing so that all robots steer in the correct direction as this is the biggest difference between steering methods.
% \begin{itemize}
%   \item Because Idris and Argo have different steering, use a param to specify which steering and program chooses weather to reverse steer or not.
%   \item Different launch file for different robots
% \end{itemize}

% Overall - With other ways that were worse so not chosen
% Tools chosen/chosen for me
% Languages
% Detailed node designs/interactions
% Messages
% Additional scripts for ease of use
% Polish

\chapter{Design after iterations and improvements}

Due to using a form of the agile methodology, each sprint iterated and improved parts of the design. This chapter goes into depth on what these changes and improvements were and why they needed to happen. Some of these changes will be referred to in the implementation chapter.

\section{Overall architecture changes}
\paragraph{}
The file format for storing the LaserScan distance data wasn't originally decided. During a sprint a couple of weeks in when the recording node was being made, a decision needed to be made for what format the values should be stored in. The data structure the file needs to store is an array of floats for each snapshot and therefore the format could be quite simple. XML and json were both considered and both rejected due to usually being used for more complex structures even though they both have tools and libraries to parse the data into the program. A much simpler format such as spreadsheets tsv (tab spaced values) or csv (comma spaced values) were considered due to being lightweight and simple to use. In the end csv format is used due to its simplicity of only needing to split the values by comma to retrieve the data and its readability from a human perspective.

\section{Detailed node designs and interaction changes}
Here is the updated detailed node designed as of the last sprint of the project. These nodes and interactions can be visualised in a diagram as shown in figure TODO.
\subsection{Recording node}
\paragraph{}
The method in which snapshots get taken was changed from relying on user input to being taken automatically depending on the distance travelled. The distance gets calculated when movement messages get received by this node and the distance since the last snapshot is calculated and then added to the distance counter. When the distance counter gets above a set threshold, a snapshot gets taken and stored in a new line in the csv file. This method repeats until the robot has driven the full path manually. When the robot is at the end of its route and the recording can be stopped, one last snapshot is taken before closing the file to make sure the end position is accurate.

This improvement made the whole system more automated as it is no longer waiting on user input. It also made the model more consistent with a set distance between snapshots and should therefore help with guiding the robot along the correct path.

\subsection{ICP compare node}
\paragraph{}
A useful metric on deciding if the match fit well enough or not is the "goodness" factor in which MRPT's ICP library generates for every match. This goodness factor represents how well the source matches the model and if this value is low then it is an indicator that the wrong model is being used. This goodness factor is therefore published with the movement data so the automatic node can use it and react accordingly by changing state if this value gets too low.

\subsection{Automatic driving node}
\paragraph{}
There were multiple issues with the original automatic driving node plan, one of which was deciding how to convert the directional transformation into robotic movement. There are multiple ways to work out where the robot needs to drive such as using pre-existing path planning algorithms or by putting weights on the y and angle values so that both eventually line up. A path planning library is built into ROS as a \cite{ROS-Path} but this way seemed unnecessary complex for this system and is usually used for checkpoints that are far away and may have obstacles in the way. I therefore went with a reactionary based system to simplify the design.

\paragraph{}
With the driving being a reaction based system, the x position is used to determine the speed the robot should be going and when to change state. The closer the robot is from the x position, the slower it should go so there aren't many anomalies. When the x position reaches a pre-determined threshold, the state can be incremented if the current state isn't the final one. When the x position is low enough and the model is on the last state, the end condition has been reached and the robot can stop due to being at the end of its route. For steering, if the y position is out by a pre-determined threshold, it gets added to the direction in which the robot needs to steer. If the angle is out by a pre-determined threshold, it also gets added to the direction in which the robot needs to steer. Weights were put onto the y and angle values in such a way that the system priorities getting into the correct position first and then the correct alignment. Once the speed and direction have been calculated, it gets put into a cmd\_vel message and published to the robot which gets guided through the checkpoints until the end state has been reached.

\section{Message changes}
\subsection{ICP direction message}
\paragraph{}
As mentioned in the ICP compare node, a goodness factor was included into the ICP\_direction message as another float and used for recognising when the source and model had a bad match. The updated message format can be seen in figure TODO.

\section{User configuration changes}
\paragraph{}
Due to the recording node improving on when snapshots are taken, no user input is needed when manually driving the robot along the path for the first time. This is a good improvement because it stops the user from having to multi-task with driving the robot and taking snapshots at the same time. It also streamlines the whole setup of the system as it simplifies what the user has to do before the automated part of the system can run.

\chapter{Software Design, Implementation and Testing}

This could be one chapter or a few chapters. It should define and discuss the software that is developed to support the research that is being conducted. For example, if your research involves running experiments, what software are you creating to support that work? What functionality is required? What design will be used? What implementation issues are there and what testing is used? 

Even though a research project is investigating specific research questions, it is still necessary for you to discuss the software that you develop. Research has a habit of generating bits of software that can exist for several years and need future modification. Therefore you need to be able to discuss the technical issues as well as the research approach. 

\section{Design}
You should concentrate on the more important aspects of the design. It is essential that an overview is presented before going into detail. As well as describing the design adopted it must also explain what other designs were considered and why they were rejected.

The design should describe what you expected to do, and might also explain areas that you had to revise after some investigation.

Typically, for an object-oriented design, the discussion will focus on the choice of objects and classes and the allocation of methods to classes. The use made of reusable components should be described and their source referenced. Particularly important decisions concerning data structures usually affect the architecture of a system and so should be described here.

How much material you include on detailed design and implementation will depend very much on the nature of the project. It should not be padded out. Think about the significant aspects of your system. For example, describe the design of the user interface if it is a critical aspect of your system, or provide detail about methods and data structures that are not trivial. Do not spend time on long lists of trivial items and repetitive descriptions. If in doubt about what is appropriate, speak to your supervisor.
 
You should also identify any support tools that you used. You should discuss your choice of implementation tools - programming language, compilers, database management system, program development environment, etc.

Some example sub-sections may be as follows, but the specific sections are for you to define. 

\subsection{Overall Architecture}

\subsection{Some detailed design}

\subsubsection{Even more detail}

\subsection{User Interface}

\subsection{Other relevant sections}

\section{Implementation}

This section should discuss issues you encountered as you tried to implement your experiments. What were the results of running the experiments? What conclusions can you draw from these results? 

During the work, you might have found that elements of your experiments were unnecessary or overly complex; perhaps third-party libraries were available that simplified some of the functions that you intended to implement. If things were easier in some areas, then how did you adapt your project to take account of your findings?

It is more likely that things were more complex than you first thought. In particular, were there any problems or difficulties that you found during implementation that you had to address? Did such problems simply delay you or were they more significant? 

If you had multiple experiments to run, it may be sensible to discuss each experiment in separate sections. 

\section{Testing}
Detailed descriptions of every test case are definitely not what is required in this section; the place for detailed lists of tests cases is in an appendix. In this section, it is more important to show that you adopted a sensible strategy that was, in principle, capable of testing the system adequately even if you did not have the time to test the system fully. 

Provide information in the body of your report and the appendix to explain the testing that has been performed. How does this testing address the requirements and design for the project?

How comprehensive is the testing within the constraints of the project?  Are you testing the normal working behaviour? Are you testing the exceptional behaviour, e.g. error conditions? Are you testing security issues if they are relevant for your project?

Have you tested your system on ``real users''? For example, if your system is supposed to solve a problem for a business, then it would be appropriate to present your approach to involve the users in the testing process and to record the results that you obtained. Depending on the level of detail, it is likely that you would put any detailed results in an appendix. 

Whilst testing with ``real users'' can be useful, don't see it as a way to shortcut detailed testing of your own. Think about issues discussed in the lectures about until testing, integration testing, etc. User testing without sensible testing of your own is not a useful activity.

The following sections indicate some areas you might include. Other sections may be more appropriate to your project. 

\subsection{Overall Approach to Testing}

\subsection{Automated Testing}

\subsubsection{Unit Tests}

\subsubsection{User Interface Testing}

\subsubsection{Stress Testing}

\subsubsection{Other types of testing}

\subsection{Integration Testing}

\subsection{User Testing}

\chapter{Results and Conclusions}

This section should discuss issues you encountered as you tried to implement your experiments. What were the results of running the experiments? What conclusions can you draw from these results? What graphs or other information have you assessed regarding your experiments? Discuss those.

During the work, you might have found that elements of your experiments were unnecessary or overly complex; perhaps third-party libraries were available that simplified some of the functions that you intended to implement. If things were easier in some areas, then how did you adapt your project to take account of your findings?

It is more likely that things were more complex than you first thought. In particular, were there any problems or difficulties that you found during implementation that you had to address? Did such problems simply delay you or were they more significant?

If you had multiple experiments to run, it may be sensible to discuss each experiment in separate sections.

\chapter{Evaluation}

% Examiners expect to find a section addressing questions such as:
\section{Requirements}
\paragraph{}
The requirements for this system were correctly identified at the beginning. The system needed to be able to automatically drive into the van and in doing so, a specific set of patterns can be detected with use of the LIDAR.

\section{Design}
\paragraph{}
The design was thought out and planned accordingly. Some changes were needed to the overall design in the amount of nodes and the node interactions.

\paragraph{}
The major changes that need to be made to make this system better if started again would be to use the ICP model matching algorithm as a service that can be called on rather than having it continuously running and sending the resulting transformation. By using it as a service, the whole system can be controlled by the automated driving node which would result in a hybrid control system which could choose if the robot needed to act upon updated information or continue with the path originally planned.

\paragraph{}
The hardest part of this project was translating the result from the ICP algorithm into movement which guided the robot along the correct path while recovering from any errors that might make it move slightly off course. An easier way of doing this would be to use pre-existing path planning tools to plan ahead so the system could be more robust.

\paragraph{}
The recording node was well thought out from the start and needed minimal changes to work as intended. When it was realised that the best way to record the data would be at specific distances rather than user defined due to the ease of use and reliability, it improved the overall workflow of the system drastically.

\section{Use of tools}
The tools that were used were suitable for this project however, a few extra tools could have been used to make the system more reliable and to help debug the issues that occurred.

\subsection{ROS path planning}
\paragraph{}
Possible tools that would have helped include using the ROS navigation stack to plan a path to the target location. This would have saved a lot of time in trying to get the robot steering to match the direction that the ICP algorithm was returning. It also would have resulted in a likely better coping system with misalignment and bad positioning due to turning the steering into a deliberative system which can react to errors rather than a fully reactive system.

\subsection{Visualisation of the data}
\paragraph{}
A tool that would have also helped would have been some sort of tool to visualise the data that the recording node took. This would have been a helpful debugging tool when figuring out certain behaviours instead of looking at the raw numbers. A tool such as rvis\cite{RVIS}, which is built into ROS, could help with such visualisations as it can take the data produced by the laserscan in real time and visualise it against the model.

\section{Testing}
\paragraph{}
The testing for this system has been thorough. None of the testing is automated but this is difficult to accomplish when working with a fully reactive robotic system. Tests were always ran multiple times to make sure the result is always reliable.

\section{Project aims}
\paragraph{}
This project set out to get Aberystwyth universities large robots to be able to automatically load themselves into a van. The final part of testing proves that the system that has been created fulfils this aim. There are however some limitations which this system has which need to be known before using it.

\paragraph{}
These limitations are that the robot must be mostly aligned with the ramp before starting, the robot must be mostly in the right position before starting and the robot must not be too far away from the ramp before starting. If all three of these limitations are met then the system will work.


% \begin{itemize}
%    \item Were the requirements correctly identified?
%    \begin{itemize}
%      \item Yes, Patern matching and moving in relation to that
%    \end{itemize}
%    \item Were the design decisions correct?
%    \begin{itemize}
%      \item Yes, it works but could have been better
%    \end{itemize}
%    \item Could a more suitable set of tools have been chosen?
%    \begin{itemize}
%      \item Maybe ROS navigation stack be used for more precice movement.
%    \end{itemize}
%    \item How well did the software meet the needs of those who were expecting to use it?
%    \begin{itemize}
%      \item TODO: Edit this after running on the real robot
%    \end{itemize}
%    \item How well were any other project aims achieved?
%    \begin{itemize}
%      \item I learned a lot more abour ROS and how to process sensordata.
%    \end{itemize}
%    \item If you were starting again, what would you do differently?
%    \begin{itemize}
%      \item Spend more time looking up movement strategies as that was the troubling part.
%      \item Use the automated movement to control everything and the ICP node as a service rather than ICP sending messages whenever it can and automated movement only partly in control, This is for easier searching.
%    \end{itemize}
% \end{itemize}
%
% Other questions can be addressed as appropriate for a project.
%
% The questions are an indication of issues you should consider. They are not intended as a specification of a list of sections.
%
% The evaluation is regarded as an important part of the project report; it should demonstrate that you are capable not only of carrying out a piece of work but also of thinking critically about how you did it and how you might have done it better. This is seen as an important part of an honours degree.
%
% There will be good things in the work and aspects of the work that could be improved. As you write this section, identify and discuss the parts of the work that went well and also consider ways in which the work could be improved.
%
% In the latter stages of the module, we will discuss the evaluation. That will probably be around week 9, although that differs each year.

% add any additional chapters here

%TC:ignore
\setemptyheader

\nocite{*} % include everything from the bibliography, irrespective of whether it has been referenced.

% the following line is included so that the bibliography is also shown in the table of contents. There is the possibility that this is added to the previous page for the bibliography. To address this, a newline is added so that it appears on the first page for the bibliography.
\addcontentsline{toc}{chapter}{Annotated Bibliography} % Adds References to contents page

%
% example of including an annotated bibliography. The current style is an author date one. If you want to change, comment out the line and uncomment the subsequent line. You should also modify the packages included at the top (see the notes earlier in the file) and then trash your aux files and re-run.
%\bibliographystyle{StylesAndReferences/authordate2annot}
\bibliographystyle{StylesAndReferences/IEEEannotU}
\renewcommand{\bibname}{Annotated Bibliography}

\bibliography{StylesAndReferences/references} % References file


\setemptyheader

\addcontentsline{toc}{chapter}{Appendices}
\chapter*{Appendices}
% The appendices are for additional content that is useful to support the discussion in the report. It is material that is not necessarily needed in the body of the report, but its inclusion in the appendices makes it easy to access.
%
% For example, if you have developed a Design Specification document as part of a plan-driven approach for the project, then it would be appropriate to include that document as an appendix. In the body of your report you would highlight the most interesting aspects of the design, referring your reader to the full specification for further detail.
%
% If you have taken an agile approach to developing the project, then you may be less likely to have developed a full requirements specification. Perhaps you use stories to keep track of the functionality and the 'future conversations'. It might not be relevant to include all of those in the body of your report. Instead, you might include those in an appendix.
%
% There is a balance to be struck between what is relevant to include in the body of your report and whether additional supporting evidence is appropriate in the appendices. Speak to your supervisor or the module coordinator if you have questions about this.
%
\pagebreak

% start the appendix - sets up different numbering
\fancypagestyle{plain}{%
%\fancyhf{} % clear all header and footer fields
\fancyhead[L]{Appendix\ \thechapter}
\fancyhead[R]{\leftmark}}

\appendix
\fancyhead[L]{Appendix\ \thechapter}
\fancyhead[R]{\leftmark}
\fancyhead[C]{}
\fancyfoot[C]{\thepage}
\renewcommand{\headrulewidth}{0.4pt}
\renewcommand{\chaptermark}[1]{\markboth{#1}{}}

\fancyhead[L]{Appendix\ \thechapter}
\fancyhead[R]{\leftmark}
\fancyfoot[C]{{\thepage} of \pageref{LastPage}}

% include any appendices here
\chapter{Third-Party Code and Libraries}

% If you have made use of any third party code or software libraries, i.e. any code that you have not designed and written yourself, then you must include this appendix. 
%
% As has been said in lectures, it is acceptable and likely that you will make use of third-party code and software libraries. If third party code or libraries are used, your work will build on that to produce notable new work. The key requirement is that we understand what your original work is and what work is based on that of other people.
%
% Therefore, you need to clearly state what you have used and where the original material can be found. Also, if you have made any changes to the original versions, you must explain what you have changed.
%
% The following is an example of what you might say.
%
% Apache POI library - The project has been used to read and write Microsoft Excel files (XLS) as part of the interaction with the client's existing system for processing data. Version 3.10-FINAL was used. The library is open source and it is available from the Apache Software Foundation
% \cite{apache_poi}. The library is released using the Apache License
% \cite{apache_license}. This library was used without modification.
%
% Include as many declarations as appropriate for your work. The specific wording is less important than the fact that you are declaring the relevant work.

% \usepackage{pdfpages}
\chapter{Ethics Submission}

% This appendix includes a copy of the ethics submission for the project. After you have completed your Ethics submission, you will receive a PDF with a summary of the comments. That document should be embedded in this report, either as images, an embedded PDF or as copied text. The content should also include the Ethics Application Number that you receive.

\includepdf[pages=-]{ethics.pdf}

% \chapter{Code Examples}

For some projects, it might be relevant to include some code extracts in an appendix. You are not expected to put all of your code here - the correct place for all of your code is in the technical submission that is made in addition to the Project Report. However, if there are some notable aspects of the code that you discuss, including that in an appendix might be useful to make it easier for your readers to access.

As a general guide, if you are discussing short extracts of code then you are advised to include such code in the body of the report. If there is a longer extract that is relevant, then you might include it as shown in the following section.

Only include code in the appendix if that code is discussed and referred to in the body of the report.

\section{Random Number Generator}

The Bayes Durham Shuffle ensures that the psuedo random numbers used in the simulation are further shuffled, ensuring minimal correlation between subsequent random outputs \cite{NumericalRecipes}.

\begin{verbatim}
 #define IM1 2147483563
 #define IM2 2147483399
 #define AM (1.0/IM1)
 #define IMM1 (IM1-1)
 #define IA1 40014
 #define IA2 40692
 #define IQ1 53668
 #define IQ2 52774
 #define IR1 12211
 #define IR2 3791
 #define NTAB 32
 #define NDIV (1+IMM1/NTAB)
 #define EPS 1.2e-7
 #define RNMX (1.0 - EPS)

 double ran2(long *idum)
 {
   /*---------------------------------------------------*/
   /* Minimum Standard Random Number Generator          */
   /* Taken from Numerical recipies in C                */
   /* Based on Park and Miller with Bays Durham Shuffle */
   /* Coupled Schrage methods for extra periodicity     */
   /* Always call with negative number to initialise    */
   /*---------------------------------------------------*/

   int j;
   long k;
   static long idum2=123456789;
   static long iy=0;
   static long iv[NTAB];
   double temp;

   if (*idum <=0)
   {
     if (-(*idum) < 1)
     {
       *idum = 1;
     }else
     {
       *idum = -(*idum);
     }
     idum2=(*idum);
     for (j=NTAB+7;j>=0;j--)
     {
       k = (*idum)/IQ1;
       *idum = IA1 *(*idum-k*IQ1) - IR1*k;
       if (*idum < 0)
       {
         *idum += IM1;
       }
       if (j < NTAB)
       {
         iv[j] = *idum;
       }
     }
     iy = iv[0];
   }
   k = (*idum)/IQ1;
   *idum = IA1*(*idum-k*IQ1) - IR1*k;
   if (*idum < 0)
   {
     *idum += IM1;
   }
   k = (idum2)/IQ2;
   idum2 = IA2*(idum2-k*IQ2) - IR2*k;
   if (idum2 < 0)
   {
     idum2 += IM2;
   }
   j = iy/NDIV;
   iy=iv[j] - idum2;
   iv[j] = *idum;
   if (iy < 1)
   {
     iy += IMM1;
   }
   if ((temp=AM*iy) > RNMX)
   {
     return RNMX;
   }else
   {
     return temp;
   }
 }

\end{verbatim}


\fancypagestyle{plain}{%
   \fancyhead{} %[C]{Annotated Bibliography}
   \fancyfoot[C]{{\thepage} of \pageref{LastPage}} % except the center
   \renewcommand{\headrulewidth}{0pt}
   \renewcommand{\footrulewidth}{0pt}
}

%TC:endignore

\end{document}
